\section[Geometric Comparison and Subcalibrations]{Geometric Comparison and Subcalibrations across Singular Sets}
\label{sec:preliminaries}

All ambient metrics in this section are Euclidean. For background on
minimal surfaces, functions of bounded variation, and minimizing
boundaries, see \cite{Giusti84}.

\subsection{Stationarity and stability on the regular part}
\label{sec:prelim-stability}

A smooth submanifold is \emph{minimal} when its mean curvature vector
vanishes. For singular objects we use integral currents, with the
boundary convention \(\partial T(\eta)=T(d\eta)\), mass measure
\(\|T\|\), and restriction \(T\llcorner A\) \cite{Federer69}.
The associated varifold is \emph{stationary} when
\begin{equation}\label{eq:stationarity-definition}
 \delta|S|(Y)=\int\diver_{T_xS}Y\,d\|S\|(x)=0,
 \qquad Y\in C_c^1(\Omega\setminus\operatorname{spt}\partial S;\R^N),
\end{equation}
where \(T_xS\) is the approximate tangent plane and its divergence
is the trace of the ambient derivative
\cite[Section~16]{Simon83}.
We write \(S_{\reg}\) for the regular part
and \(\Sing S\) for the remaining support, away from the boundary.

For a smooth minimal \(M\), let \(A\) be its second fundamental form
and \(\nabla^\perp\) its normal connection.  The second variation in a
compactly supported normal direction \(V\) is
\begin{equation}\label{eq:normal-second-variation}
 Q_M^\perp(V)=\int_M\left(|\nabla^\perp V|^2
       -\sum_{i,j=1}^n\langle A(e_i,e_j),V\rangle^2\right),
\end{equation}
where \(e_1,\ldots,e_n\) is a local orthonormal tangent frame.
With the Laplacian convention \(\int u\Delta u=-\int|\nabla u|^2\),
the normal Jacobi operator is
\begin{equation}\label{eq:normal-Jacobi-definition}
 J_M^\perp V=\Delta_M^\perp V+
       \sum_{i,j}\langle A(e_i,e_j),V\rangle A(e_i,e_j),
 \qquad Q_M^\perp(V)=-\int_M\langle V,J_M^\perp V\rangle.
\end{equation}
Stability means \(Q_M^\perp(V)\geq0\) for every such \(V\).
For a stationary singular current the tests are compactly supported
smooth normal fields on \(S_{\reg}\), away from its boundary;
see \cite{Simons68} and
\cite[Lemma~9.7 and Remark~9.8]{Simon83}.
Each test has compact support away from the singular set, but no common
positive distance is prescribed for the whole test class.
Extension across it uses intrinsic energy cutoffs, whereas divergence
comparison below uses ambient \(L^1\) cutoffs. For cones, strict
stability means the scale-invariant gap defined below.

\subsection{Area minimization and the strict notions for cones}
\label{sec:prelim-minimization}

An integral current \(S\) is \emph{locally mass minimizing} if
\[
 \mathbf M_U(S)\leq\mathbf M_U(S+\partial R)
\]
whenever \(U\Subset\Omega\) and \(R\) is an integral current compactly
supported in \(U\) of one higher dimension.  Here
\(\mathbf M_U(S)=\|S\|(U)\).  For multiplicity-one
submanifolds we also call this \emph{area minimization}.
Such minimizers are stationary and stable, by first and second variation.

A measurable set \(E\subset\Omega\) has \emph{locally finite perimeter}
if \(\one_E\in BV_{\rm loc}(\Omega)\): its distributional derivative
is a vector-valued Radon measure of locally finite mass. For
\(U\Subset\Omega\),
\begin{equation}\label{eq:finite-perimeter-definition}
 \Per(E;U)=|D\one_E|(U),\qquad
 D\one_E=-\nu_E\,\cH^{N-1}\llcorner\partial^*E,
\end{equation}
where \(\partial^*E\) is the reduced boundary and \(\nu_E\) its
measure-theoretic exterior unit normal; see \cite{Maggi12}.

The current \(\llbracket E\rrbracket\) denotes integration over \(E\)
with the standard orientation, and
\(\mathbf M_U(\partial\llbracket E\rrbracket)=\Per(E;U)\).
Its boundary has multiplicity one on \(\partial^*E\), with support
\(\overline{\partial^*E}\cap\Omega\). For polynomial hypersurfaces we take
\(C=\partial\llbracket\{F<0\}\rrbracket\).
Its support equals the full zero set when both signs occur near every
zero, in particular for nonzero harmonic \(F\);
see \cref{rem:polynomial-finite-perimeter}.

A set \(E\) is \emph{locally perimeter minimizing} if
\(\Per(E;U)\leq\Per(H;U)\) for every locally finite-perimeter set
\(H\) with \(H\triangle E\Subset U\Subset\Omega\).
Local perimeter minimization of \(E\) is equivalent to local mass
minimization of \(\partial\llbracket E\rrbracket\).  This is a direct
consequence of the codimension-one decomposition theorem
\cite[Theorem~27.6 and Corollary~27.8]{Simon83}; see also
\cite[Proposition~2.15]{DeLellis15}.

For the cone statements let \(C\) be a stationary multiplicity-one
integral cone of real dimension \(n\geq2\), with \(\partial C=0\).
Put \(r=|x|\) and \(C_1=C\llcorner B_1\).
We also use \(C\) for its support. By slicing and homogeneity,
\(\partial C_1=\langle C,r,1\rangle\) is the oriented spherical slice
on \(\partial B_1\) \cite[Lemma~28.5(2)]{Simon83}.

\begin{definition}[Minimality and stability, with their strict forms]
\label{def:four-variational-notions}
The cone is \emph{area minimizing} if \(C_1\) has least mass among
integral \(n\)-currents with boundary \(\partial C_1\).  It is
\emph{strictly area minimizing} in the vertex-avoidance sense if it is
area minimizing and there is \(\Theta>0\) such that
\begin{equation}\label{eq:strict-minimality-definition}
 \mathbf M(T)-\mathbf M(C_1)\geq\Theta\eps^n
 \quad\text{whenever}\quad
 \partial T=\partial C_1,\quad
 \operatorname{spt}T\cap B_\eps=\varnothing,\quad 0<\eps<1.
\end{equation}
The competitors are integral currents of finite mass and arbitrary
integer multiplicity in the ambient Euclidean space.
For \(C=\partial\llbracket E\rrbracket\) with \(E\)
conical, area minimization is equivalent to local perimeter minimization.

For a two-sided hypersurface cone and \(V=u\nu_C\),
\eqref{eq:normal-second-variation} becomes
\begin{equation}\label{eq:Q-definition}
 Q_C(u)=\int_{C_{\reg}}\bigl(|\nabla_Cu|^2-|A_C|^2u^2\bigr),
 \qquad u\in C_c^\infty(C_{\reg}\setminus\{0\}).
\end{equation}
The cone is \emph{stable} if \(Q_C(u)\geq0\) for every such \(u\), and
\emph{strictly stable} if there is \(\delta>0\) such that
\begin{equation}\label{eq:strict-stability-definition}
 Q_C(u)\geq\delta\int_{C_{\reg}}r^{-2}u^2
 \qquad\text{for every }u\in C_c^\infty(C_{\reg}\setminus\{0\}).
\end{equation}
In arbitrary codimension the same definition uses \(Q_C^\perp(V)\)
and \(|V|^2\) in place of \(Q_C(u)\) and \(u^2\).
\end{definition}

The vertex-avoidance convention is that of Hardt--Simon
\cite[Section~3, (3.1)]{HardtSimon85}, stated for integral currents in
\cite[p.~210]{Lin87}. These strict notions enter constructions with
prescribed conical singularities \cite{CaffarelliHardtSimon84},
metric deformations \cite{WangDeformations20}, and Simon's Liouville
theorem \cite[equations~(1.4), (1.12) and Corollary~1.13]{Simon21}.

Strict stability does not imply area minimization,
as \cref{prop:iso-stability,thm:classical-iso-classification} illustrate;
\cref{prop:complex-zero-gap} gives calibrated area-minimizing cones
of real codimension two without a positive Hardy gap.
For a singular link
\(\Lambda=C\cap\Sph^{N-1}\), angular tests are compactly supported in
\(\Lambda_{\reg}\); the corresponding spectral formulation is recorded
in \eqref{eq:exact-Hardy-gap}.

\subsection{Calibrations and the Stokes comparison}
\label{sec:prelim-calibrations}

A \emph{calibration} is a closed differential \(n\)-form whose
comass, the supremum of its absolute value on orthonormal
\(n\)-frames, is at most one. It calibrates an integral current \(S\) if
\(\omega(\tau_S)=1\) at \(\|S\|\)-almost every point, where \(\tau_S\)
is the oriented unit tangent \(n\)-vector
\cite{HarveyLawson82a}.

If \(\omega\) calibrates \(S\) and \(T-S=\partial R\) with
\(\operatorname{spt}R\Subset U\), Stokes' theorem gives
\begin{equation}\label{eq:calibration-Stokes-comparison}
 \mathbf M_U(T)-\mathbf M_U(S)
 \geq (T-S)(\omega)=\partial R(\omega)=R(d\omega)=0.
\end{equation}
The common contributions outside \(U\) cancel; a cutoff equal to one
near \(\operatorname{spt}R\) makes the evaluation of the difference
compactly supported.  Thus a calibrated current is mass minimizing.

\subsection{Hypersurface forms and vector-field subcalibrations}
\label{sec:prelim-subcalibrations}

In codimension one, the identification \(\omega_X=\iota_XdV\) gives
\begin{equation}\label{eq:form-vector-equivalence}
 \|\omega_X\|_*=|X|,\qquad
 d\omega_X=(\diver X)dV,\qquad
 \omega_X|_{\partial E}=(X\cdot\nu_E)\,d\mu_{\partial E}.
\end{equation}
Evaluation on an oriented orthonormal frame and the coordinate
formula for divergence give these identities. A hypersurface
calibration is equivalently a field with \(|X|\leq1\),
\(\diver X=0\), and \(X\cdot\nu_E=1\) on the boundary.

\begin{definition}[Two-sided subcalibration, smooth formulation]
\label{def:subcalibration}
Let \(E\subset\Omega\subset\R^N\) have smooth boundary.  A
\(C^1\) field \(X\) is a two-sided subcalibration for \(E\) if
\begin{equation}\label{eq:smooth-subcalibration-signs}
 |X|\leq1,\qquad X\cdot\nu_E=1\text{ on }\partial E,\qquad
 \diver X\leq0\text{ in }E,\quad
 \diver X\geq0\text{ in }\Omega\setminus E.
\end{equation}
The equivalent form conditions follow from \eqref{eq:form-vector-equivalence}.
\end{definition}

These phase signs give the comparison used by De~Philippis and Paolini
\cite{DePhilippisPaolini}. Applying the Stokes argument above to a smooth
competitor \(H\) with \(H\triangle E\Subset U\Subset\Omega\) gives
\begin{align}
 \Per(H;U)-\Per(E;U)
 &\geq\int_U(\one_H-\one_E)\diver X\,dx\notag\\
 &=\int_{H\triangle E}q_X\,dx\geq0,
 \qquad q_X=(1-2\one_E)\diver X.
 \label{eq:smooth-subcalibration-Stokes}
\end{align}
The boundary flux is saturated for \(E\) and bounded by the perimeter
for \(H\). Below we extend this comparison to singular and
phase-dependent fields.

\begin{remark}[Historical background on quadratic cones]
\label{rem:subcalibration-history}
Fleming's regularity theorem in $\R^3$ \cite{Fleming62} and
Almgren's theorem in $\R^4$ \cite{Almgren66} exclude nonflat
area-minimizing hypercones in those dimensions.
Simons established stability of the balanced quadratic cones in
\(\R^{2m}\), \(m\geq4\), \cite{Simons68}; Bombieri--De~Giorgi--Giusti
proved their area minimality \cite[Sections~II--III]{BDGG69}.
The homogeneous barriers in \cite[Section~IV, p.~260]{BDGG69}
occur in their construction of entire minimal graphs, separately from
the cone-minimization proof. Lawson reduced the equivariant Plateau
problem to weighted orbit spaces and constructed exact calibrating
one-forms through explicit homogeneous trial functions
\cite[Section~5, pp.~244--247]{Lawson72}.
In his 1973 Berkeley thesis, Sim\~oes completed the remaining
quadratic cases: the cone over the minimal product
$S^1\times S^5$ is not area minimizing, whereas the cone over
the minimal product $S^2\times S^4$ is area minimizing
\cite{Simoes73}. Here the factors have the radii making their
product minimal in the unit sphere.
Miranda and collaborators developed new proofs through the
nonparametric minimal surface equation. In particular,
Massari--Miranda used graphical subsolutions and cylinder limits
\cite{MassariMiranda83}; the quartic function
\(G=(|x|^4-|y|^4)/4\) already appears in this approach
\cite[pp.~146--147]{Miranda08}.
Benarros--Miranda gave two one-sided homogeneous functions
\cite[Section~2, pp.~46--47]{BenarrosMiranda95}.
After orienting both normalized gradient fields consistently, their
construction gives potentials \(G_\pm=a_\pm F\) with positive,
phase-dependent multipliers.
De~Philippis--Paolini combined the earlier approaches in a short
subcalibration proof, using divergence inequalities directly in the
ambient space \cite[pp.~234--235]{DePhilippisPaolini}.
De~Philippis--Maggi used the exact perimeter-deficit identity
\cite[Proposition~4.1, equation~(4.5)]{DePhilippisMaggi14}
to obtain quantitative stability estimates for Lawson cones.
Liu completed their six exceptional parameter pairs using different
subcalibrations \cite[Theorems~1--2 and Section~2.1]{Liu19},
explicitly acknowledging the inspiration of Benarros--Miranda
\cite[p.~1002]{Liu19}.
\end{remark}

\subsection{Comparison across removable exceptional sets}
\label{sec:comparison}

The Stokes comparison \eqref{eq:smooth-subcalibration-Stokes}
extends across locally \(\cH^{N-1}\)-null singular sets,
also for different phase potentials.

\subsection*{Boundary flux away from an exceptional set}

For a locally finite-perimeter set \(E\), a relatively closed set
\(Z\), and \(X\in C^1(\Omega\setminus Z;\R^N)\), the boundary
condition is \(X\cdot\nu_E=1\) almost everywhere on
\(\partial^*E\setminus Z\). Extending \(\eta X\) by zero for
\(\eta\in C_c^1(\Omega\setminus Z)\), Gauss--Green gives
\[
 \int_E\diver(\eta X)\,dx
 =\int_{\partial^*E\setminus Z}\eta\,X\cdot\nu_E\,d\cH^{N-1}.
\]
To extend this calculation across \(Z\), we use cutoffs below.
For singular subcalibration fields in a Riemannian metric, see
\cite[Fact~2.3 and Section~7]{LiuSimon26}.

\subsection*{Cutoffs near the singular set}

\begin{lemma}[Cutoffs near sets of zero Hausdorff measure]
\label{lem:null-set-cutoffs}
Let \(Z\subset\Omega\subset\R^N\) be relatively closed and satisfy
\begin{equation}\label{eq:H-null}
 \cH^{N-1}(Z\cap K)=0
 \quad\text{for every compact }K\Subset\Omega.
\end{equation}
For every compact \(K\Subset U\Subset\Omega\), every open set
\(O\) with \(Z\cap K\subset O\Subset U\), and every \(\eps>0\),
there is \(\psi\in C_c^1(O)\) such that \(0\leq\psi\leq1\),
\(\psi=1\) on a neighborhood of \(Z\cap K\), and
\begin{equation}\label{eq:null-set-cutoff}
 \int_U\bigl(\psi+|\nabla\psi|\bigr)<\eps.
\end{equation}
\end{lemma}

\begin{proof}
Cover the compact set \(Z\cap K\) by finitely many balls
\(B_{r_i}(x_i)\) whose doubled balls lie in \(O\), with
\(r_i<1\) and \(\sum_i r_i^{N-1}<\eps/(2C_N)\).
Choose smooth cutoffs \(\psi_i\) supported in the doubled balls,
equal to one on the covering balls, with \(0\leq\psi_i\leq1\)
and \(|\nabla\psi_i|\leq C/r_i\).
Then \(\psi=1-\prod_i(1-\psi_i)\) equals one near \(Z\cap K\),
and satisfies \(\psi\leq\sum_i\psi_i\) and
\(|\nabla\psi|\leq\sum_i|\nabla\psi_i|\).  Hence
\[
 \int_U\bigl(\psi+|\nabla\psi|\bigr)
 \leq C_N\sum_i(r_i^N+r_i^{N-1})<\eps.\qedhere
\]
\end{proof}

We extend the quantitative subcalibration identity of De Philippis--Maggi
\cite[Proposition~4.1, equation~(4.5)]{DePhilippisMaggi14}, also used in
\cite[Section~2, equation~(3)]{Liu19}, to bounded fields that are \(C^1\)
outside a relatively closed, locally \(\cH^{N-1}\)-null set.

\Needspace{7\baselineskip}
\begin{theorem}[Two-sided comparison]
\label{thm:removable-subcalibration}
Let \(E\subset\Omega\) have locally finite perimeter, and let
\(Z\subset\Omega\) be relatively closed with
\(\cH^{N-1}(Z\cap K)=0\) for every compact \(K\Subset\Omega\).
Suppose
\(X\in C^1(\Omega\setminus Z;\R^N)\) satisfies
\begin{enumerate}[label=\textup{(\roman*)}]
 \item \(|X|\leq1\);
 \item \(X\cdot\nu_E=1\) for \(\cH^{N-1}\)-almost every point of
 \(\partial^*E\setminus Z\);
 \item \(\diver X\leq0\) almost everywhere in \(E\setminus Z\), and
 \(\diver X\geq0\) almost everywhere in \((\Omega\setminus E)\setminus Z\).
\end{enumerate}
Define \(q_X=(1-2\one_E)\diver X\geq0\) off \(Z\) and set it to zero on
\(Z\).  Then \(q_X\in L^1_{\rm loc}(\Omega)\).  For every
locally finite-perimeter competitor \(H\) with \(H\triangle E\Subset U\Subset\Omega\),
\begin{align}
 \Per(H;U)-\Per(E;U)
 &=\int_{\partial^*H\cap U}(1-X\cdot\nu_H)\,d\cH^{N-1}
   +\int_{H\triangle E}q_X\,dx.\label{eq:quantitative-comparison-identity}
\end{align}
The values of \(X\) on \(Z\) do not affect the boundary integral.
In particular,
\begin{equation}\label{eq:quantitative-comparison}
 \Per(H;U)-\Per(E;U)\geq\int_{H\triangle E}q_X\,dx.
\end{equation}
Thus \(E\) is locally perimeter minimizing and
\(\partial\llbracket E\rrbracket\) is locally mass minimizing.
If \(q_X>0\) almost everywhere, equality in perimeter comparison forces
\(H=E\) up to a null set.  When \(|X|=1\), the first integral in
\eqref{eq:quantitative-comparison-identity} is
\(\frac12\int_{\partial^*H\cap U}|\nu_H-X|^2\,d\cH^{N-1}\).
\end{theorem}

\begin{proof}
Choose \(K_0\Subset K_1\Subset U\) containing the support of
\(\one_H-\one_E\) in \(K_0\), and put
\(\mu=|D\one_E|+|D\one_H|\).
The null-set assumption and \eqref{eq:finite-perimeter-definition}
give \(\mu(Z\cap K_1)=0\).
Take open neighborhoods \(O_j\Subset U\) of \(Z\cap K_1\) with
\(\mu(O_j)\to0\), and use \cref{lem:null-set-cutoffs} to choose
\(\psi_j\in C_c^1(O_j)\) with
\(\int_U(\psi_j+|\nabla\psi_j|)\to0\).
Pass to a subsequence so that \(\psi_j\to0\) almost everywhere, and set
\(\eta_j=1-\psi_j\).  Fix \(\chi\in C_c^1(U)\) equal to one near
\(K_0\), with support in \(K_1\).
The field \(\chi\eta_jX\), extended by zero near \(Z\), is an admissible
\(C^1\) field.  The distributional derivative identity
\eqref{eq:finite-perimeter-definition}, applied to
\(\one_H-\one_E\), gives
\begin{align}
 \int_{H\triangle E}\eta_jq_X
 &=\int_{\partial^*H}\chi\eta_j X\cdot\nu_H
   -\int_{\partial^*E}\chi\eta_j
   -\int_U(\one_H-\one_E)X\cdot\nabla\eta_j.
 \label{eq:BV-cutoff-comparison}
\end{align}
We have used \(\chi=1\) on the difference set and
\((\one_H-\one_E)\diver X=q_X\) there.  The final integral is bounded in
absolute value by \(\int_U|\nabla\psi_j|\).
The boundary integrals converge because their cutoff errors are at most
\(\mu(O_j)\). Thus the right-hand side of
\eqref{eq:BV-cutoff-comparison} has a finite limit. Since
\(0\leq\eta_jq_X\to q_X\) almost everywhere, Fatou's lemma gives
\[
 \int_{H\triangle E}q_X
 \leq\liminf_j\int_{H\triangle E}\eta_jq_X<\infty.
\]
Now \(0\leq\eta_jq_X\leq q_X\), so dominated convergence gives
equality in the limit of \eqref{eq:BV-cutoff-comparison}.
Since \(H=E\) outside \(K_0\), the boundary fluxes cancel there.
Consequently
\[
 \int_{H\triangle E}q_X
 =\int_{\partial^*H\cap U}X\cdot\nu_H-\Per(E;U),
\]
which proves \eqref{eq:quantitative-comparison-identity}.
For almost every ball \(B\Subset\Omega\), both
\(H=E\setminus B\) and \(H=E\cup B\) have locally finite perimeter.
The integrability just proved for these two comparisons gives integrability
of \(q_X\) on the two parts of \(B\), hence local integrability in
\(\Omega\).  The current conclusion follows from the equivalence in
\cref{sec:prelim-minimization}; the other conclusions follow from
nonnegativity.
\end{proof}

\begin{remark}[Semialgebraic singular sets]
A closed semialgebraic set of dimension at most \(N-2\) satisfies
\eqref{eq:H-null}: its finite smooth stratification
\cite[Sections~2.3 and~2.8]{BochnakCosteRoy98} has strata of zero
\(\cH^{N-1}\)-measure.
\end{remark}

\begin{remark}[Finite perimeter of polynomial sublevel sets]
\label{rem:polynomial-finite-perimeter}
Let \(F\) be a nonzero real polynomial of degree \(d\), and put
\(E=\{F<0\}\) and \(Q_L=(-L,L)^N\).
On each coordinate line, the restriction of \(F\) is either zero
or a nonzero univariate polynomial of degree at most \(d\).
Thus the restriction of \(\one_E\) has one-dimensional variation
at most \(d\). For \(\varphi\in C_c^1(Q_L)\), Fubini and
one-dimensional integration by parts give
\[
 \left|\int_{Q_L}\one_E\,\partial_i\varphi\,dx\right|
 \leq d(2L)^{N-1}\|\varphi\|_\infty.
\]
Each distributional derivative \(D_i\one_E\) is therefore a finite
measure on \(Q_L\), and
\[
 \Per(E;Q_L)\leq\sum_{i=1}^N|D_i\one_E|(Q_L)
 \leq Nd(2L)^{N-1}.
\]
In particular, polynomial sublevel sets have locally finite perimeter.
This also applies to \(\{G<0\}\) whenever \(G\) has the same sign
almost everywhere as a polynomial.

We also recall that \(\{F=0\}\) has Lebesgue measure zero.
Write \(F\) as a polynomial in one coordinate. Outside the zero
set of a nonzero coefficient polynomial, its fibers are finite.
Induction on \(N\) and Fubini prove the assertion.

If both signs of \(F\) occur in every neighborhood of each zero, then
\[
 \partial E=\operatorname{spt}\partial\llbracket E\rrbracket
 =\{F=0\}.
\]
Indeed, both phases have positive volume in every ball centered at
a zero. The perimeter in such a ball cannot vanish, since otherwise
\(\one_E\) would be constant almost everywhere there. Away from
the zero set, \(\one_E\) is locally constant.
For nonzero harmonic \(F\), the strong maximum principle gives this
sign condition.
\end{remark}

\begin{corollary}[Gradient subcalibration criterion]
\label{cor:gradient-subcalibration}
Let \(G\in C^2(\Omega\setminus Z)\cap C^0(\Omega)\), assume that
\(E=\{G<0\}\) has locally finite perimeter, and suppose that \(Z\) contains
the critical set of \(G\) and the singular part of \(\{G=0\}\).
If \(Z\) is relatively closed, \(\cH^{N-1}(Z\cap K)=0\) for every
compact \(K\Subset\Omega\), and
\[
 G\cL(G)\geq0\qquad\text{on }\Omega\setminus Z,
\]
then \(E\) is locally perimeter minimizing.
\end{corollary}

\begin{proof}
Use \(X=\nabla G/|\nabla G|\) and
\eqref{eq:div-normalized-gradient}.  On \(G<0\) its divergence is
nonpositive, and on \(G>0\) it is nonnegative.  Along the regular zero set it
is the exterior unit normal.
\end{proof}

We call \(G\) a \emph{sub-calibration potential} for \(E=\{G<0\}\)
when its normalized gradient satisfies \cref{def:subcalibration}, or
\cref{thm:removable-subcalibration} when an exceptional set is present.
\Cref{def:homogeneous-subcalibration-potential} specifies the homogeneous
potentials used for polynomial defining equations.

\paragraph{Applying the criterion to homogeneous potentials.}
In the algebraic applications, \(G=aF\) has positive degree
\(\ell\), and \(a\) is smooth and positive outside a closed conical
exceptional set \(Z\) containing the singular zeros of \(F\).
Once \(G\) extends continuously with the signs of \(F\), its negative
phase has locally finite perimeter by
\cref{rem:polynomial-finite-perimeter}. Moreover, outside \(Z\),
\[
 F\neq0\ \Longrightarrow\
 \langle x,\nabla G\rangle=\ell G\neq0,
 \qquad
 F=0\ \Longrightarrow\ \nabla G=a\nabla F\neq0.
\]
It remains to check the null-set condition on \(Z\) and the global
sign \(G\cL(G)\geq0\). The exceptional sets and their codimensions
are collected in \cref{app:exceptional-set-summary}.

\subsection*{One-sided comparison and two different fields}

A set \(E\) is \emph{inward minimizing} if perimeter comparison holds
for competitors contained in \(E\), and \emph{outward minimizing} if
it holds for competitors containing \(E\).  All changes are compactly
supported. These are the subminimality and complementary subminimality
of \cite[Definition~1.1]{DePhilippisPaolini}. The paired comparison
retains both divergence remainders.

\begin{theorem}[One-sided and paired comparison]
\label{thm:paired-subcalibration}
Let \(E\subset\Omega\) have locally finite perimeter, and let
\(Z\subset\Omega\) be relatively closed with
\(\cH^{N-1}(Z\cap K)=0\) for every compact \(K\Subset\Omega\).
Suppose \(X_-,X_+\in C^1(\Omega\setminus Z;\R^N)\) satisfy
\[
 |X_\pm|\leq1,\qquad X_\pm\cdot\nu_E=1
 \quad\cH^{N-1}\text{-a.e. on }\partial^*E\setminus Z.
\]
Assume only the respective phase signs
\[
 \diver X_-\leq0\text{ in }E\setminus Z,\qquad
 \diver X_+\geq0\text{ in }(\Omega\setminus E)\setminus Z.
\]
Define, with zero values on \(Z\),
\[
 q_-=-\one_E\diver X_-,\qquad
 q_+=\one_{\Omega\setminus E}\diver X_+.
\]
Then \(q_-,q_+\in L^1_{\rm loc}(\Omega)\), and the following
conclusions hold for changes compactly supported in \(U\Subset\Omega\).
\begin{enumerate}[label=\textup{(\roman*)}]
 \item The assumptions on \(X_-\) alone imply
 \[
  \Per(H;U)-\Per(E;U)\geq\int_{E\setminus H}q_-
  \qquad(H\subset E).
 \]
 The assumptions on \(X_+\) alone imply the corresponding estimate
 \[
  \Per(H;U)-\Per(E;U)\geq\int_{H\setminus E}q_+
  \qquad(H\supset E).
 \]
 \item Together they imply, for every locally finite-perimeter \(H\),
 \begin{equation}\label{eq:paired-weighted-comparison}
  \Per(H;U)-\Per(E;U)
  \geq\int_{E\setminus H}q_-+\int_{H\setminus E}q_+.
 \end{equation}
 Thus \(E\) is locally perimeter minimizing, and its oriented boundary
 is locally mass minimizing among integral currents.  If each weight
 is positive almost everywhere on its phase, equality implies
 \(H=E\) up to a null set.
\end{enumerate}
\end{theorem}

\begin{proof}
For \(H\subset E\), the cutoff proof of
\cref{thm:removable-subcalibration} uses the divergence sign only on
\(E\setminus H\). With \(X_-\), it gives
\[
 \Per(H;U)-\Per(E;U)
 =\int_{\partial^*H\cap U}(1-X_-\cdot\nu_H)
  +\int_{E\setminus H}q_-.
\]
The analogous identity with \(X_+\) holds for \(H\supset E\).
Their nonnegative boundary terms prove \textup{(i)}; taking
\(H=E\setminus B\) and \(H=E\cup B\) for almost every relatively
compact ball proves local integrability of the respective weights.

For an unrestricted competitor, apply these inequalities to
\(H\cap E\) and \(H\cup E\).  The perimeter inequality for
unions and intersections \cite{Maggi12} gives
\[
 \Per(H\cap E;U)+\Per(H\cup E;U)
 \leq\Per(H;U)+\Per(E;U).
\]
Adding the two one-sided estimates proves
\eqref{eq:paired-weighted-comparison}.  The current conclusion follows
from the equivalence in \cref{sec:prelim-minimization}, and positivity
gives the equality assertion.
\end{proof}

\begin{remark}[Fields defined on their respective sides]
\label{rem:phase-local-fields}
When \(E\) is open and \(\partial E\setminus Z\) is smooth, it is
enough to construct \(X_-\) on an open neighborhood of
\(\overline E\setminus Z\), and \(X_+\) on an open neighborhood of
\(\overline{\Omega\setminus E}\setminus Z\), relative to
\(\Omega\setminus Z\).  A smooth cutoff equal to one on the relevant
closed phase extends each bounded field by zero to the unused side.
This preserves its norm, boundary values, and the divergence on its
own phase.  In a conical setting the cutoff may be chosen on the unit
sphere, so homogeneity is preserved.  No matching of the two potentials
is required; their normalized gradients both equal \(\nu_E\) on the
regular boundary.

Liu's subcalibrations for the six exceptional Lawson-cone parameter pairs
\cite[Sections~2.1 and~3]{Liu19} use powers of different block radii,
and sometimes different homogeneous degrees, on the two sides.
The positive multipliers \(G_\pm=a_\pm F\) give the common normal
\(\nu_E\) along the regular cone, so
\cref{thm:paired-subcalibration} applies.

Phase-dependent multipliers also appear in Liu's application of
Simon's analytic estimates to hypersurfaces with prescribed closed
singular sets \cite[Theorem~1 and Definition~1.3]{LiuSimon26}.
\end{remark}

\subsection*{Strict minimality and distance-weighted comparison}

\begin{theorem}[Strict minimality from two-sided divergence integrals]
\label{thm:strict-minimality-flux}
Let \(E\subset\R^{n+1}\), \(n\geq2\), be conical, and suppose the hypotheses
of \cref{thm:removable-subcalibration} hold globally with a conical
exceptional set \(Z\) and a zero-homogeneous field \(X\).
Suppose
\begin{equation}\label{eq:two-phase-flux}
 A_-:=\int_{B_1\cap E}q_X\,dx>0,
 \qquad
 A_+:=\int_{B_1\setminus E}q_X\,dx>0.
\end{equation}
Then \(C=\partial\llbracket E\rrbracket\) is strictly area minimizing in
\cref{def:four-variational-notions}, with
\begin{equation}\label{eq:strict-minimality-constant}
 \Theta=\min\{A_-,A_+\}.
\end{equation}
In particular the assertion permits nonisolated singularities.
\end{theorem}

\begin{proof}
Both integrals are finite by \cref{thm:removable-subcalibration}.
Homogeneity gives \(q_X(rx)=r^{-1}q_X(x)\), so their values over
\(B_\eps\) are \(\eps^nA_-\) and \(\eps^nA_+\).
If a set competitor \(H\) has no perimeter in the connected ball
\(B_\eps\), then \(\one_H\) is constant there.  Accordingly,
\(H\triangle E\) contains either \(B_\eps\cap E\) or
\(B_\eps\setminus E\), up to a null set.
Equation \eqref{eq:quantitative-comparison} proves the desired mass gap
for set-boundary competitors.

For an arbitrary integral-current competitor, consider the nearest-point
projection onto the closed unit ball,
\[
 \pi(x)=\frac{x}{\max\{1,|x|\}}.
\]
This map is \(1\)-Lipschitz, fixes \(\partial C_1\), and maps
\(\R^{n+1}\setminus B_\eps\) into itself.
Although \(\pi\) is not proper, the integral representation of the
finite-mass current \(T\) defines \(\pi_\#T\), since the
\(n\)-Jacobian of \(\pi\) is bounded by one.
The pushforward mass estimate and the identity
\(\partial(\pi_\#T)=\pi_\#(\partial T)\)
give
\[
 \mathbf M(\pi_\#T)\leq\mathbf M(T),
 \qquad \partial(\pi_\#T)=\partial C_1;
\]
see \cite[Remark~26.21(1), Lemma~26.25, and
Remark~27.2(3)]{Simon83}.  For finite-mass currents with noncompact
support, these statements follow from the compactly supported case by
cutoffs: the additional boundary term is bounded by
\(O(L^{-1})\mathbf M(T)\) for a cutoff at radius \(L\), and tends to
zero.  The same projection is used in
\cite[proof of Proposition~1.2]{DeLellisSpadaroSpolaor17}.
We may therefore replace \(T\) by \(\pi_\#T\) and assume
\(\operatorname{spt}T\subset\overline B_1\).

Since \(\partial(T-C_1)=0\), coning \(T-C_1\) to the origin gives an
integral \((n+1)\)-current \(R\) with
\(\partial R=T-C_1\) and support in \(\overline B_1\)
\cite[equation~(26.26) and Remark~27.2(2)--(3)]{Simon83}.
A top-dimensional integral current is integration against an
integer-valued multiplicity.  Since its boundary has finite mass, this
multiplicity has bounded variation
\cite[Remark~26.28]{Simon83}.  Thus
\[
 R=\llbracket k\rrbracket,
 \qquad k\in BV_c(\R^{n+1};\mathbb Z),
 \qquad v:=\one_E+k\in BV_{\mathrm{loc}}(\R^{n+1};\mathbb Z).
\]
We have \(\partial\llbracket v\rrbracket=T+C-C_1\).
In \(B_\eps\), the right-hand side vanishes because \(T\) avoids
this ball and \(C=C_1\) there.  The Constancy Theorem
\cite[Theorem~26.27]{Simon83}, applied on the connected
open set \(B_\eps\), therefore gives \(v=m\) almost everywhere
there for some \(m\in\mathbb Z\).

Set \(H=\{v\geq1\}\). The integer-level decomposition
\cite[Theorem~27.6]{Simon83} gives, for every relatively compact
open set \(U\subset\R^{n+1}\),
\[
 |Dv|(U)=\sum_{j\in\mathbb Z}
       \Per(\{v\geq j\};U)
 \geq\Per(H;U).
\]
Thus \(H\) has locally finite perimeter. Moreover,
\(H=E\) outside \(\overline B_1\), and \(\one_H\) is constant
in \(B_\eps\).
Since \(\|C\|(\partial B_1)=0\), the mass measures of \(T\) and
\(C-C_1\) are mutually singular.  Hence, for \(\rho>1\),
\[
 \Per(H;B_\rho)\leq |Dv|(B_\rho)
 =\mathbf M(T)+\Per(E;B_\rho)-\mathbf M(C_1).
\]
Apply the set comparison in \(B_\rho\) and cancel the common exterior
perimeter.  This yields
\eqref{eq:strict-minimality-definition} with
\eqref{eq:strict-minimality-constant}.
\end{proof}


\begin{corollary}[Strict mass gap from paired fields]
\label{cor:paired-strict-minimality}
Suppose \cref{thm:paired-subcalibration} holds globally for a conical
set \(E\subset\R^{n+1}\), \(n\geq2\), a conical exceptional set \(Z\),
and fields that are zero-homogeneous on their respective phases.
If
\[
 A_-:=\int_{B_1\cap E}q_->0,\qquad
 A_+:=\int_{B_1\setminus E}q_+>0,
\]
then every integral-current competitor in
\eqref{eq:strict-minimality-definition} satisfies
\[
 \mathbf M(T)-\mathbf M(C_1)
 \geq\min\{A_-,A_+\}\eps^n.
\]
\end{corollary}

\begin{proof}
The degree-\(-1\) weights are locally integrable by
\cref{thm:paired-subcalibration}. A set competitor avoiding \(B_\eps\)
fills or empties it, so \eqref{eq:paired-weighted-comparison} gives the
gap. The current-to-set reduction in the proof of
\cref{thm:strict-minimality-flux} preserves ball avoidance and does
not increase the mass excess.
\end{proof}

\begin{corollary}[Polynomial potentials with nonzero divergence]
\label{cor:polynomial-strict-minimality}
Let \(N\geq3\).  Suppose a homogeneous polynomial \(G\) satisfies the global hypotheses of
\cref{cor:gradient-subcalibration}.  If both \(\{G<0\}\) and \(\{G>0\}\)
are nonempty and \(\cL(G)\not\equiv0\), then
\(\partial\llbracket\{G<0\}\rrbracket\) is strictly area minimizing.
\end{corollary}

\begin{proof}
The normalized gradient is zero-homogeneous.  Its critical set is
conical and satisfies \eqref{eq:H-null}.  Off this set,
\[
 q_X=\frac{|\cL(G)|}{|\nabla G|^3}.
\]
A nonzero polynomial has a Lebesgue-null zero set
by \cref{rem:polynomial-finite-perimeter}. Thus \(q_X>0\) almost everywhere on both open
phases.  Equation \eqref{eq:two-phase-flux} follows, and
\cref{thm:strict-minimality-flux} applies.
\end{proof}

For nonpolynomial potentials, verify positive flux in both phases and
apply \cref{thm:strict-minimality-flux}, as for the radical potentials in
\cref{ex:Riedler-correct-power,sec:complex-determinants}.

For regular area-minimizing hypercones, Niu \cite[Theorem~3]{Niu26}
characterizes simultaneous strict stability and strict minimality by
\eqref{eq:distance-weighted-comparison}. We prove sufficiency across
the singular sets of \cref{thm:removable-subcalibration}, using the
quantitative subcalibration method of De Philippis--Maggi
\cite[Section~4]{DePhilippisMaggi14} and Liu
\cite[Section~2.2]{Liu19}.

\begin{proposition}[Distance-weighted sufficient criterion]
\label{prop:distance-weighted-comparison}
Under the geometric and comparison hypotheses of
\cref{thm:strict-minimality-flux}, suppose the cone is two-sided and minimal on its regular part and, for some \(c>0\),
\begin{equation}\label{eq:linear-distance-divergence}
 q_X(x)\geq c\,\frac{\dist(x,C)}{|x|^2}
 \quad\text{for almost every }x\notin C.
\end{equation}
Assume each phase contains points at positive distance from \(C\).
Then the cone is strictly area minimizing and strictly stable, and
\begin{equation}\label{eq:distance-weighted-comparison}
 \Per(H;U)-\Per(E;U)
 \geq c\int_{H\triangle E}\frac{\dist(x,C)}{|x|^2}\,dx.
\end{equation}
More precisely, \(Q_C(u)\geq c\int_Cr^{-2}u^2\) for the regular-part
test class in \cref{def:four-variational-notions}.
\end{proposition}

\begin{proof}
The weighted comparison follows from
\eqref{eq:quantitative-comparison}. The two phase integrals are
positive, so \cref{thm:strict-minimality-flux} gives strict minimality.
For a smooth compactly supported regular-part function \(u\), perturb the
boundary by the normal graph \(tu\) in a fixed regular tubular neighborhood.
Minimality and tubular coordinates give
\begin{align*}
 \Per(E_{tu};U)-\Per(E;U)&=\tfrac12t^2Q_C(u)+o(t^2),\\
 \int_{E_{tu}\triangle E}\frac{\dist(x,C)}{|x|^2}\,dx
 &=\tfrac12t^2\int_Cr^{-2}u^2+o(t^2).
\end{align*}
Indeed, on each normal segment the leading integral is
\(r^{-2}\int_0^{|tu|}s\,ds\); its tubular Jacobian and the radial weight
contribute only higher-order terms on this compact support.
Divide \eqref{eq:distance-weighted-comparison} by \(t^2/2\) and let
\(t\to0\).
\end{proof}

\subsection{Algebraic tangent cones and local instability}
\label{sec:tangent-instability-transfer}

Let $C=\{F=0\}$ for a nonzero polynomial $F$, and let $p\in C$.
Writing $F(p+z)=H_m(z)+\cdots+H_d(z)$ with $H_m\not\equiv0$
homogeneous of degree $m$, we have
\[
 r^{-m}F(p+rz)\longrightarrow H_m(z)
 \quad\text{in }C^\infty_{\mathrm{loc}}\quad(r\downarrow0).
\]
The implicit function theorem gives smooth graphical convergence of
$(C-p)/r$ near compact subsets of $\{H_m=0,\ \nabla H_m\ne0\}$.
If $H_m$ changes sign near every zero, then $\{H_m=0\}$ is the
unique local Hausdorff tangent: uniform convergence gives one inclusion,
and sign change gives the other. This condition holds when $F$ is harmonic
or $H_m$ is an indefinite quadratic form. Thus an indefinite Hessian at
a critical point gives the quadratic tangent
$\{z:\langle\Hess F(p)z,z\rangle=0\}$.
For homogeneous $F$ and $p\ne0$, Euler's identity gives $D_pH_m=0$,
so this leading zero set splits off the line $\R p$.
The cubic Hessian criterion in \cref{prop:cubic-general-singular-spectrum}
identifies these quadratic factors.

\begin{lemma}[Instability under a regular tangent limit]
\label{lem:tangent-instability-transfer}
Let $C$ be an $n$-dimensional hypersurface, minimal on its two-sided
regular part. Suppose that $C_j=(C-p)/r_j$, $r_j\downarrow0$,
converges locally graphically in $C^2$ on compact regular subsets of
$\Sigma=Q\times\R^k$, where $Q$ is a two-sided minimal cone.
If $Q$ is unstable on its regular part, then $C$ admits a compactly
supported regular-part variation with negative second variation
in every neighborhood of $p$.
\end{lemma}

\begin{proof}
Write $Q_\Sigma(\phi)=\int_\Sigma
(|\nabla_\Sigma\phi|^2-|A_\Sigma|^2\phi^2)$.
Choose $\phi\in C_c^\infty(Q_{\reg})$ with $Q_Q(\phi)<0$.
For $k>0$, take a nonzero $\chi\in C_c^\infty(\R^k)$ and
set $\chi_T(w)=\chi(w/T)$. The product metric gives
\[
 Q_{Q\times\R^k}(\phi\chi_T)
 =T^k\left(\int_{\R^k}\chi^2\right)Q_Q(\phi)
 +T^{k-2}\left(\int_{\R^k}|\nabla\chi|^2\right)
             \left(\int_Q\phi^2\right)<0
\]
for sufficiently large $T$; for $k=0$ use $\phi$ itself.
Fix the resulting test function $\psi$ and a relatively compact regular
neighborhood of its support. Transport $\psi$ by the graphical
parametrizations to $\psi_j$ on $C_j$, extending by zero.
The induced metrics and second fundamental forms converge, so
$Q_{C_j}(\psi_j)\to Q_\Sigma(\psi)<0$.
For $u_j(x)=\psi_j((x-p)/r_j)$,
\[
 Q_C(u_j)=r_j^{\,n-2}Q_{C_j}(\psi_j)<0.
\]
Its support lies in the regular part and shrinks to $p$.
\end{proof}

A leading form $a(|u|^2-|v|^2)$ in orthogonal coordinates,
$a\ne0$, with $u,v\in\R^s$ and $s=2,3$, gives an unstable
quadratic factor: $|A|^2=(2s-2)/\rho^2$, where
$\rho^2=|u|^2+|v|^2$, and the sharp radial Hardy constant is
$(2s-3)^2/4<2s-2$.
For example, the minimal cubic $F=\operatorname{Re}(z_1z_2z_3)$
on $\mathbb C^3$ has leading term $\operatorname{Re}(z_2z_3)$
at $p=(1,0,0)$, giving the tangent
$\R^2\times\{(u,v)\in\R^2\oplus\R^2:|u|=|v|\}$.
The lemma proves local instability at $p$; see
\cref{sec:tangents} for the Clifford and Hurwitz cubic models.
